\subsection{最小化最坏剩余半径的模型}

如果希望避免某些反馈带来很差的定位结果，可以改用最坏情形准则。对于每个检测点 \(q\)，先考察所有可能反馈，将其中最大的定位损失记为 \(J_{\mathrm W}(q)\)。再比较不同检测点，选择这个最大损失最小的位置，即

\begin{equation}
 \min_{q\in Q}\quad J_{\mathrm W}(q)=\sup_{z\in F(q)}L_z(q).
\label{eq:q2-pair-worst-objective}
\end{equation}

这个模型不需要假设各反馈出现的概率。只要某种反馈符合题意、没有被已有观测排除，就必须考虑其后果，而不能因为它很少出现就忽略它。因此，该方法给出的是不利情况下仍能达到的定位效果。

对于同一个检测点，平均损失不会超过最大的可能损失，即 \(J_{\mathrm E}(q)\le J_{\mathrm W}(q)\)。但两者最小值所在的位置可能不同：一个位置可能平均效果更好，另一个位置则可能在最不利反馈下表现更好。

